%!TeX encoding = UTF-8
%!TeX program = xelatex
\documentclass[notheorems, aspectratio=54]{beamer}
% aspectratio: 1610, 149, 54, 43(default), 32

\usepackage{latexsym}
\usepackage{amsmath,amssymb}
\usepackage{mathtools}
\usepackage{color,xcolor}
\usepackage{graphicx}
\usepackage{algorithm}
\usepackage{amsthm}
% font hỗ trợ tiếng Việt (biên dịch bằng xelatex)
\usepackage{fontspec}
\setmainfont{DejaVuSerif.ttf}[Path=/usr/share/fonts/truetype/dejavu/]
\setsansfont{DejaVuSans.ttf}[Path=/usr/share/fonts/truetype/dejavu/]
\usepackage{animate} % insert gif

\usepackage{lipsum} % To generate test text 
\usepackage{ulem} % 下划线，波浪线

\usepackage{listings} % display code on slides; don't forget [fragile] option after \begin{frame}

% ----------------------------------------------
% tikx
\usepackage{framed}
\usepackage{tikz}
\usepackage{pgf}
\usetikzlibrary{calc,trees,positioning,arrows,chains,shapes.geometric,%
    decorations.pathreplacing,decorations.pathmorphing,shapes,%
    matrix,shapes.symbols}
\pgfmathsetseed{1} % To have predictable results
% Define a background layer, in which the parchment shape is drawn
\pgfdeclarelayer{background}
\pgfsetlayers{background,main}

% define styles for the normal border and the torn border
\tikzset{
  normal border/.style={orange!30!black!10, decorate, 
     decoration={random steps, segment length=2.5cm, amplitude=.7mm}},
  torn border/.style={orange!30!black!5, decorate, 
     decoration={random steps, segment length=.5cm, amplitude=1.7mm}}}

% Macro to draw the shape behind the text, when it fits completly in the
% page
\def\parchmentframe#1{
\tikz{
  \node[inner sep=2em] (A) {#1};  % Draw the text of the node
  \begin{pgfonlayer}{background}  % Draw the shape behind
  \fill[normal border] 
        (A.south east) -- (A.south west) -- 
        (A.north west) -- (A.north east) -- cycle;
  \end{pgfonlayer}}}

% Macro to draw the shape, when the text will continue in next page
\def\parchmentframetop#1{
\tikz{
  \node[inner sep=2em] (A) {#1};    % Draw the text of the node
  \begin{pgfonlayer}{background}    
  \fill[normal border]              % Draw the ``complete shape'' behind
        (A.south east) -- (A.south west) -- 
        (A.north west) -- (A.north east) -- cycle;
  \fill[torn border]                % Add the torn lower border
        ($(A.south east)-(0,.2)$) -- ($(A.south west)-(0,.2)$) -- 
        ($(A.south west)+(0,.2)$) -- ($(A.south east)+(0,.2)$) -- cycle;
  \end{pgfonlayer}}}

% Macro to draw the shape, when the text continues from previous page
\def\parchmentframebottom#1{
\tikz{
  \node[inner sep=2em] (A) {#1};   % Draw the text of the node
  \begin{pgfonlayer}{background}   
  \fill[normal border]             % Draw the ``complete shape'' behind
        (A.south east) -- (A.south west) -- 
        (A.north west) -- (A.north east) -- cycle;
  \fill[torn border]               % Add the torn upper border
        ($(A.north east)-(0,.2)$) -- ($(A.north west)-(0,.2)$) -- 
        ($(A.north west)+(0,.2)$) -- ($(A.north east)+(0,.2)$) -- cycle;
  \end{pgfonlayer}}}

% Macro to draw the shape, when both the text continues from previous page
% and it will continue in next page
\def\parchmentframemiddle#1{
\tikz{
  \node[inner sep=2em] (A) {#1};   % Draw the text of the node
  \begin{pgfonlayer}{background}   
  \fill[normal border]             % Draw the ``complete shape'' behind
        (A.south east) -- (A.south west) -- 
        (A.north west) -- (A.north east) -- cycle;
  \fill[torn border]               % Add the torn lower border
        ($(A.south east)-(0,.2)$) -- ($(A.south west)-(0,.2)$) -- 
        ($(A.south west)+(0,.2)$) -- ($(A.south east)+(0,.2)$) -- cycle;
  \fill[torn border]               % Add the torn upper border
        ($(A.north east)-(0,.2)$) -- ($(A.north west)-(0,.2)$) -- 
        ($(A.north west)+(0,.2)$) -- ($(A.north east)+(0,.2)$) -- cycle;
  \end{pgfonlayer}}}

% Define the environment which puts the frame
% In this case, the environment also accepts an argument with an optional
% title (which defaults to ``Example'', which is typeset in a box overlaid
% on the top border
\newenvironment{parchment}[1][Example]{%
  \def\FrameCommand{\parchmentframe}%
  \def\FirstFrameCommand{\parchmentframetop}%
  \def\LastFrameCommand{\parchmentframebottom}%
  \def\MidFrameCommand{\parchmentframemiddle}%
  \vskip\baselineskip
  \MakeFramed {\FrameRestore}
  \noindent\tikz\node[inner sep=1ex, draw=black!20,fill=white, 
          anchor=west, overlay] at (0em, 2em) {\sffamily#1};\par}%
{\endMakeFramed}

% ----------------------------------------------

\mode<presentation>{
    \usetheme{CambridgeUS}
    % Boadilla CambridgeUS
    % default Antibes Berlin Copenhagen
    % Madrid Montpelier Ilmenau Malmoe
    % Berkeley Singapore Warsaw
    \usecolortheme{beaver}
    % beetle, beaver, orchid, whale, dolphin
    \useoutertheme{infolines}
    % infolines miniframes shadow sidebar smoothbars smoothtree split tree
    \useinnertheme{circles}
    % circles, rectanges, rounded, inmargin
}
% 设置 block 颜色
\setbeamercolor{block title}{bg=red!30,fg=white}

\newcommand{\reditem}[1]{\setbeamercolor{item}{fg=red}\item #1}

% 缩放公式大小
\newcommand*{\Scale}[2][4]{\scalebox{#1}{\ensuremath{#2}}}

% 解决 font warning
\renewcommand\textbullet{\ensuremath{\bullet}}

% ---------------------------------------------------------------------
% flow chart
\tikzset{
    >=stealth',
    punktchain/.style={
        rectangle, 
        rounded corners, 
        % fill=black!10,
        draw=white, very thick,
        text width=6em,
        minimum height=2em, 
        text centered, 
        on chain
    },
    largepunktchain/.style={
        rectangle,
        rounded corners,
        draw=white, very thick,
        text width=10em,
        minimum height=2em,
        on chain
    },
    line/.style={draw, thick, <-},
    element/.style={
        tape,
        top color=white,
        bottom color=blue!50!black!60!,
        minimum width=6em,
        draw=blue!40!black!90, very thick,
        text width=6em, 
        minimum height=2em, 
        text centered, 
        on chain
    },
    every join/.style={->, thick,shorten >=1pt},
    decoration={brace},
    tuborg/.style={decorate},
    tubnode/.style={midway, right=2pt},
    font={\fontsize{10pt}{12}\selectfont},
}
% ---------------------------------------------------------------------

% code setting
\lstset{
    language=C++,
    basicstyle=\ttfamily\footnotesize,
    keywordstyle=\color{red},
    breaklines=true,
    xleftmargin=2em,
    numbers=left,
    numberstyle=\color[RGB]{222,155,81},
    frame=leftline,
    tabsize=4,
    breakatwhitespace=false,
    showspaces=false,               
    showstringspaces=false,
    showtabs=false,
    morekeywords={Str, Num, List},
}

% ---------------------------------------------------------------------

%% preamble
\title[Smart locker — Cầu Giấy]{Tối ưu vị trí đặt trạm dịch vụ (smart locker)\\quận Cầu Giấy cũ}
\author{INT3501E — Service Science}
\institute[VNU]{Học phần INT3501E}

% -------------------------------------------------------------

\begin{document}

%% title frame
\begin{frame}
    \titlepage
\end{frame}

\section{Bài toán}

\begin{frame}
    \frametitle{Bài toán}
    Chọn \textbf{tập vị trí đặt trạm dịch vụ (smart locker) nhỏ nhất} trên địa bàn
    quận Cầu Giấy cũ (3 phường: Cầu Giấy, Nghĩa Đô, Yên Hòa; $\approx 12{,}46$ km$^2$),
    thỏa mãn:
    \begin{itemize}
        \item Chỉ đặt tại \textbf{điểm ứng viên hợp lệ} lấy từ OpenStreetMap;
              ứng viên nằm trong khu dân cư (H) hoặc khu vực cấm (C) bị loại.
        \item Mỗi người dân phát sinh \textbf{2 đơn vị nhu cầu/ngày} — tổng
              $\approx 588.166$ đơn vị/ngày từ 1.544 điểm cầu (WorldPop, lưới $\sim$100 m).
        \item Mỗi điểm cầu được gán \textbf{trọn vẹn} cho \textbf{trạm mở gần nhất};
              điểm cầu được coi là phục vụ nếu khoảng cách Manhattan đến trạm $\le R$.
        \item Tổng nhu cầu được phục vụ $\ge \alpha \cdot 100\%$ và không trạm nào
              vượt công suất ngày.
    \end{itemize}
\end{frame}

\begin{frame}
    \frametitle{Dữ liệu đầu vào}
    \begin{itemize}
        \item \textbf{Ứng viên} (OpenStreetMap): bãi đỗ xe, siêu thị, cửa hàng tiện lợi,
              bưu điện (ứng viên tòa chung cư đang bị vô hiệu hóa) — 481 điểm.
        \item \textbf{Nhu cầu} (WorldPop): 1.544 điểm cầu, tổng $\approx 588.166$ đơn vị/ngày.
        \item \textbf{Polygon}: khu dân cư (H) và khu cấm (C: quân sự, đường sắt,
              nghĩa trang, mặt nước).
        \item \textbf{Tham số trạm}: công suất ngày và chi phí đầu tư theo loại trạm.
        \item Khoảng cách: \textbf{Manhattan (L1)} trên hệ chiếu UTM EPSG:32648.
    \end{itemize}

    \begin{block}{Lọc ứng viên theo khu dân cư — phương án A (\texttt{landuse=residential})}
        \textbf{129/481} ứng viên hợp lệ: bãi đỗ xe 71, cửa hàng tiện lợi 46,
        siêu thị 9, bưu điện 3.
    \end{block}
\end{frame}

\section{Mô hình}

\begin{frame}
    \frametitle{Mô hình MILP — biến và ràng buộc}
    Biến nhị phân: $x_j = 1$ nếu mở trạm $j$; $y_{ij} = 1$ nếu điểm cầu $i$
    được gán cho trạm $j$ (chỉ xét các cặp có khoảng cách $d_{ij} \le R$).

    \begin{align*}
        \sum_{j} y_{ij} &\le 1 && \forall i &\text{(mỗi điểm cầu gán tối đa một trạm)}\\
        y_{ij} &\le x_j && \forall i,j &\text{(chỉ gán vào trạm đã mở)}\\
        x_j &\le \sum_{k \in P(i,j)} y_{ik} && \forall i,j &\text{(gán trạm gần nhất — prefix)}\\
        \sum_{i} q_i\, y_{ij} &\le Q_j\, x_j && \forall j &\text{(công suất ngày)}\\
        \sum_{i,j} q_i\, y_{ij} &\ge \alpha \sum_{i} q_i && &\text{(độ phủ nhu cầu)}
    \end{align*}

    với $q_i$: nhu cầu điểm cầu $i$; $Q_j$: công suất trạm $j$;
    $P(i,j)$: tập ứng viên gần hơn (hoặc bằng) $j$ đối với điểm cầu $i$.
\end{frame}

\begin{frame}
    \frametitle{Hàm mục tiêu hai pha}
    \begin{block}{Pha 1 — tối thiểu số trạm}
        $$\min \sum_{j} x_j$$
    \end{block}
    \begin{block}{Pha 2 — cố định số trạm tối ưu $k^*$, tối thiểu tổng chi phí đầu tư}
        $$\min \sum_{j} f_j\, x_j \quad \text{s.t.} \quad \sum_{j} x_j = k^*$$
    \end{block}
    \begin{itemize}
        \item Ưu tiên giảm số trạm trước, sau đó chọn tổ hợp trạm rẻ nhất.
        \item Cài đặt bằng \texttt{scipy.optimize.milp} (solver HiGHS),
              ma trận ràng buộc thưa.
        \item Nếu pha 2 không hội tụ trong thời gian cho phép,
              giữ nghiệm tốt nhất của pha 1.
    \end{itemize}
\end{frame}

\section{Pipeline}

\begin{frame}
    \frametitle{Pipeline xử lý}
    \begin{enumerate}
        \item \textbf{fetch\_data}: tải ranh giới, ứng viên OSM, nhu cầu WorldPop.
        \item \textbf{polygons}: lọc ứng viên theo khu H/C, tính cận trên
              độ phủ và năng lực.
        \item \textbf{distances\_l1}: ma trận khoảng cách Manhattan (EPSG:32648).
        \item \textbf{run.py}: precheck $\to$ MILP hai pha $\to$ kiểm định độc lập
              $\to$ xuất kết quả.
    \end{enumerate}
    \begin{itemize}
        \item Precheck loại sớm cấu hình bất khả thi mà không cần gọi solver.
        \item \textbf{evaluate}: gán lại nearest một cách độc lập, kiểm tra
              coverage / capacity / objective.
        \item Mỗi lần chạy ghi vào \texttt{results/<run\_id>/} riêng, không ghi đè.
        \item Kiểm thử: đối chiếu brute-force trên 40 instance ngẫu nhiên,
              pytest offline, CI trong Docker.
    \end{itemize}
\end{frame}

\section{Kết quả}

\begin{frame}
    \frametitle{Tiền kiểm khả thi (precheck)}
    Kiểm tra polygon với $R = 500$ m, $\alpha = 0{,}95$:
    \begin{center}
    \resizebox{\textwidth}{!}{%
    \begin{tabular}{ l | c c c }
        Phương án xác định H & Ứng viên hợp lệ & Cận trên độ phủ & Khả thi? \\
        \hline
        A: \texttt{landuse=residential} & 129 / 481 & 92,73\% & không (chỉ đạt khi $R = 800$ m) \\
        B: khối nhà ở (building footprints) & 184 / 481 & 97,94\% & có
    \end{tabular}}
    \end{center}
    \begin{itemize}
        \item Phương án A: cận trên độ phủ 92,73\% $<$ 95\% $\Rightarrow$
              bất khả thi ngay từ precheck khi $R = 500$ m.
        \item Toàn bộ các lần chạy thực nghiệm dùng phương án A (129 ứng viên).
    \end{itemize}
\end{frame}

\begin{frame}
    \frametitle{Kết quả thực nghiệm trên dữ liệu thật}
    \begin{center}
    \resizebox{\textwidth}{!}{%
    \begin{tabular}{ l | c c c | l }
        Cấu hình & $R$ (m) & $\alpha$ & cap & Kết quả \\
        \hline
        R500/$\alpha$0.95/cap1 & 500 & 0,95 & $\times 1$ & precheck: bất khả thi (UB độ phủ 92,73\%) \\
        R800/$\alpha$0.95/cap1 & 800 & 0,95 & $\times 1$ & bất khả thi (solver) \\
        R800/$\alpha$0.90/cap1 & 800 & 0,90 & $\times 1$ & bất khả thi (solver) \\
        R800/$\alpha$0.95/cap2 & 800 & 0,95 & $\times 2$ & \textbf{khả thi: 19 trạm}
    \end{tabular}}
    \end{center}
    \begin{itemize}
        \item cap: hệ số nhân công suất ngày của mỗi trạm.
        \item $R = 500$: precheck chặn sớm; ba cấu hình còn lại do solver
              chứng minh bất khả thi.
    \end{itemize}
\end{frame}

\begin{frame}
    \frametitle{Kết quả tốt nhất: $R = 800$ m, $\alpha = 0{,}95$, cap$\times$2}
    \begin{block}{Phương án đề xuất}
        Mở \textbf{19 trạm} — tổng chi phí đầu tư \textbf{14,5 tỷ đồng} —
        độ phủ \textbf{95,33\%} (560.688 / 588.165 đơn vị nhu cầu/ngày).
    \end{block}
    \begin{itemize}
        \item Kiểm định độc lập: \textbf{PASS} cả 5 mục — id ứng viên hợp lệ,
              gán nearest, độ phủ, công suất, nhất quán mục tiêu.
        \item Trạng thái \texttt{feasible\_time\_limit}: pha 1 tìm được nghiệm
              trong 900 s nhưng chưa kịp chứng minh tối ưu; pha 2 hết thời gian
              nên dùng nghiệm pha 1.
        \item Precheck trước khi giải: cận trên độ phủ 99,33\% $\ge$ 95\%,
              năng lực đủ.
    \end{itemize}
\end{frame}

\begin{frame}
    \frametitle{Thảo luận}
    \begin{itemize}
        \item $R = 500$ m bất khả thi về độ phủ: 129 ứng viên không phủ nổi
              nhu cầu dày — cần tăng $R$ lên 800 m hoặc thêm ứng viên
              (phương án B).
        \item cap$\times$1 bất khả thi kể cả khi hạ $\alpha = 0{,}90$: ràng buộc
              \textbf{gán trạm gần nhất} dồn nhu cầu vào một số trạm gây quá tải —
              nút thắt là công suất, không phải độ phủ.
        \item cap$\times$2 khả thi: chỉ cần nhân đôi công suất là hấp thụ
              được điểm nghẽn này.
        \item Precheck giúp phát hiện sớm bất khả thi (R500) mà không tốn
              thời gian chạy solver.
    \end{itemize}
\end{frame}

\begin{frame}
    \frametitle{Kết luận}
    \begin{itemize}
        \item Pipeline hoàn chỉnh từ dữ liệu thật (OSM + WorldPop) đến kết quả
              được kiểm định độc lập.
        \item Mô hình MILP hai pha với ràng buộc gán trạm gần nhất; cài đặt
              được đối chiếu brute-force trên 40 instance ngẫu nhiên.
        \item Kết quả tốt nhất: 19 trạm, chi phí 14,5 tỷ đồng,
              phủ 95,33\% nhu cầu.
        \item Hạn chế: khoảng cách Manhattan (chưa phải mạng đường thực tế),
              ứng viên chung cư đang tắt, pha 2 chưa hội tụ do giới hạn thời gian.
        \item Hướng tiếp theo: phương án B (building footprints), khoảng cách đi bộ,
              tăng thời gian giải cho pha 2.
    \end{itemize}
\end{frame}

\end{document}